%%%%%%%%%%%%%%%%%%%%%%%%%%%%%%%%%%%%%%%%%%%%%%%%%%%%%%%%%%%%%%%
%
% Welcome to Overleaf --- just edit your LaTeX on the left,
% and we'll compile it for you on the right. If you open the
% 'Share' menu, you can invite other users to edit at the same
% time. See www.overleaf.com/learn for more info. Enjoy!
%
%%%%%%%%%%%%%%%%%%%%%%%%%%%%%%%%%%%%%%%%%%%%%%%%%%%%%%%%%%%%%%%
\documentclass[12pt]{article}
\usepackage{iftex}
\ifPDFTeX
  \usepackage[utf8]{inputenc}
  \usepackage[T1]{fontenc}
\else
  \usepackage[no-math]{fontspec}
\fi
\usepackage{amsmath}
\usepackage{amssymb}
\usepackage{amsthm}
\usepackage{mathtools}
\usepackage[margin=2.5cm]{geometry}
\usepackage{tikz}
\usepackage{setspace}
\usepackage{quantikz}
\usepackage{fancyhdr}
\setlength{\headheight}{15pt}
\setstretch{1.25}

\usetikzlibrary{arrows,arrows.meta,automata,positioning,patterns,calc,shapes,decorations.pathreplacing}


\pagestyle{fancy}

\fancyhf{}

\fancyhead[L]{Evolución de sistemas abiertos}
\fancyhead[R]{Práctica 4}

\renewcommand{\headrulewidth}{0.4pt}
\renewcommand{\footrulewidth}{0pt}

\fancyfoot[C]{\thepage}

\title{\textbf{Seminario de Mecánica Cuántica}\\[0.3cm] \large Práctica 4}
\date{}
\author{Milton Plasencia}

\makeatletter
\AtBeginDocument{%
  \theoremstyle{definition}\@ifundefined{exercise}{\newtheorem{exercise}{Ejercicio}}{}%
}
\makeatother
\renewcommand{\proofname}{Demostración}

\begin{document}
\maketitle
\thispagestyle{fancy}

% --- Página 1 ---

\begin{exercise}
Mostremos el caso más general, el caso puro se verá como un caso particular.

Si $\rho_{AB}$ es un estado producto $\rho_{AB}=\rho_A \otimes \rho_B$ con $\rho_B$ en general estado mezcla, entonces

\begin{equation*}
     \rho_B=\sum_i p_i |\psi^B_i\rangle\langle\psi^B_i|.
\end{equation*}

En general lo que podemos hacer es escribir la evolución de $\rho_{A}$ como la traza parcial de B respecto la evolución unitaria del sistema completo, es decir,

\begin{align*}
    \rho_A' &= \text{Tr}_B\left[U_{AB}(\rho_A \otimes \rho_B)U_{AB}^{\dagger}\right] = \sum_j \langle j_B| \sum_i p_i \left[U_{AB}(\rho_A \otimes |\psi^B_i\rangle\langle\psi^B_i|)U_{AB}^{\dagger}\right]|j_B\rangle \\
    &= \sum_{i,j} \sqrt{p_i} \langle j_B|U_{AB}|\psi^B_i\rangle \rho_A \sqrt{p_i}  \langle\psi^B_i|U_{AB}^{\dagger}|j_B\rangle \\
\end{align*}

Definimos los operadores de Kraus como $E_{ij} = \sqrt{p_i} \langle j_B|U_{AB}|\psi^B_i\rangle$. Nótese que efectivamente son de Krauss pues

\begin{align*}
    \sum_{i,j} E_{ij}^{\dagger}E_{ij} &= \sum_{i,j} p_i \langle\psi^B_i|U_{AB}^{\dagger}|j_B\rangle \langle j_B|U_{AB}|\psi^B_i\rangle \\
    &= \sum_i p_i \langle\psi^B_i|U_{AB}^{\dagger}U_{AB}|\psi^B_i\rangle = \sum_i p_i \langle\psi^B_i|\mathbb{I_A}\otimes \mathbb{I_B}|\psi^B_i\rangle = \mathbb{I}_A \sum_i p_i = \mathbb{I}_A
\end{align*}

\end{exercise}



\begin{exercise}
        \textbf{a)} Para el caso de la evolución dada por un CNOT, tenemos que para el estado dado $\rho_{AB}=\rho_A \otimes|\psi^B\rangle\langle\psi^B|$, los operadores de Krauss son fácilmente computables en la base $|\pm\rangle$ de $X$, donde $X|\pm\rangle=\pm |\pm \rangle$.
En efecto, tenemos que como $U_{AB}=|0\rangle\langle 0|\otimes \mathbb{I}_B + |1\rangle\langle 1|\otimes X$ luego, si admitimos una descomposición para $|\psi_B\rangle = \alpha |+\rangle + \beta |-\rangle$, entonces

$$E_+= \langle+_B|U_{AB}|\psi_B\rangle = |0\rangle\langle 0|\otimes \langle+|\mathbb{I}_B|\psi_B\rangle + |1\rangle\langle 1| \otimes \langle-|X|\psi_B\rangle = \alpha \mathbb{I}_A $$

$$E_-= \langle-_B|U_{AB}|\psi_B\rangle = |0\rangle\langle 0|\otimes \langle-|\mathbb{I}_B|\psi_B\rangle + |1\rangle\langle 1| \otimes \langle-|X|\psi_B\rangle = \beta Z. $$

De este modo, la evolución del subsistema $\rho_A$ está dada por

$$\rho_A'=\sum_\alpha E_\alpha\rho_AE^{\dagger}= |\alpha|^2 \rho_A + |\beta|^2 Z\rho_AZ. $$

Y como son estados normalizados podemos poner $p=|\beta|^2$ y por tanto $|\alpha|^2=1-p$, obteniendo finalmente

$$\rho_A'=(1-p)\rho_A + p Z\rho_AZ. $$

Analicemos y discutamos los siguientes casos:

\textbf{Caso $p=1/2$:}

Consideremos 

\begin{equation}\label{rhoA}
  \rho_A = \begin{pmatrix}
    \rho_{00} & \rho_{01} \\
    \rho_{10} & \rho_{11}
\end{pmatrix}
\end{equation}

y hagamoslo evolucionar por el canal de Krauss, obteniendo

\begin{align*}
    \rho_A' &= \frac{1}{2}\begin{pmatrix}
        \rho_{00} & \rho_{01} \\
        \rho_{10} & \rho_{11}
    \end{pmatrix} + \frac{1}{2}\begin{pmatrix}
        \rho_{00} & -\rho_{01} \\
        -\rho_{10} & \rho_{11}
    \end{pmatrix} = \begin{pmatrix}
        \rho_{00} & 0 \\
        0 & \rho_{11}
    \end{pmatrix}.
\end{align*}

Lo que vemos es que la evolución no unitaria destruye por completo las coherencias cuánticas. Por eso esta evolución se llama \textit{decoherencia completa}.

En términos de la entropía, si suponemos que $\rho_A$ es estado puro, la variación de entropía entre $\rho_A$ y $\rho_A'$ es 

$$\Delta S = S(\rho_A') - S(\rho_A) = -\rho_{00}\log\rho_{00} - \rho_{11}\log\rho_{11} + \text{Tr}(\rho_A\log\rho_A)= .$$

\textbf{Caso $p=0$:}

Considerando nuevamente el estado \ref{rhoA}, en este caso la evolución de Krauss es trivial, pues

$$\rho_A' = \rho_A.$$

En este caso, el estado original queda inalterado, y por tanto no hay decoherencia. Este es el caso de \textit{ausencia de decoherencia}.

Como es evidente, la variación de entropía entre $\rho_A$ y $\rho_A'$ es nula. Al no haber decoherencia, no hay pérdida de información, y por tanto la entropía no varía.




\end{exercise}



\textbf{b)} Analicemos la evolución en el caso de $\rho_B$ no estado puro sino mezcla. En este caso, podemos expandir $\rho_B$ como

$$\rho_B =  \rho_{00} |0\rangle\langle 0| + \rho_{01} |0\rangle\langle 1| + \rho_{10} |1\rangle\langle 0| + \rho_{11} |1\rangle\langle 1|. $$

La evolución de $\rho_A$ está dada por

$$    \rho_A' = \text{Tr}_B\left[U_{AB}(\rho_A \otimes \rho_B)U_{AB}^{\dagger}\right] $$

Desarrollando por partes, obtenemos

\begin{align*}
   U_{AB}(\rho_A \otimes \rho_B)U_{AB}^{\dagger} &= \left(|0\rangle\langle0|\rho_A \otimes \rho_B + |1\rangle\langle1|\rho_A \otimes X\rho_B \right) (|0\rangle\langle0| \otimes \mathbb{I}_B + |1\rangle\langle1| \otimes X) \\
  &= |0\rangle\langle0|\rho_A|0\rangle\langle0| \otimes \rho_B  + |0\rangle\langle0|\rho_A|1\rangle\langle1|\otimes \rho_BX +  |1\rangle\langle1|\rho_A|0\rangle\langle0| \otimes X\rho_B + \\
  & \quad |1\rangle\langle1|\rho_A|1\rangle\langle1| \otimes X\rho_BX 
\end{align*}

y notando $\text{Tr}(\rho O)=\langle O \rangle$ y por propiedad cíclicla dentro de la traza $\text{Tr}(X\rho_A X)=\text{Tr}(\underbrace{XX}_{\mathbb{I}}\rho_A)=1$ tenemos que

\begin{align*}
  \rho_A' &=  \text{Tr}_B(U_{AB}(\rho_A \otimes \rho_B)U_{AB}^{\dagger}) \\
   &= \langle0|\rho_A|0\rangle |0\rangle\langle0|  + \langle0|\rho_A|1\rangle \langle X \rangle_B |0\rangle\langle1| + \langle1|\rho_A|0\rangle \langle X \rangle_B |1\rangle\langle0| + \langle1|\rho_A|1\rangle \langle X \rangle_B |1\rangle\langle1| \\
   &= \begin{pmatrix}
    \rho^{A}_{00} & \langle X \rangle_B \rho^{A}_{01} \\
    \langle X \rangle_B \rho^{A}_{10} & \rho^{A}_{11}
    \end{pmatrix}
\end{align*}

Si ahora ponemos $\rho_B = \frac{1}{d_B}\mathbb{I}_B$ veremos que recuperamos el caso de decoherencia completa, pues $\langle X \rangle_B = 0$ y por tanto 

$$\rho_A' = \begin{pmatrix}.
    \rho^{A}_{00} & 0 \\
    0 & \rho^{A}_{11}
    \end{pmatrix}.$$






\begin{exercise}
    \textbf{a)} Dados los operadores de Krauss $E_1 = |0\rangle\langle 0| + \sqrt{1-p}|1\rangle\langle 1|$ y $E_2 = \sqrt{p}|0\rangle\langle 1|$, podemos ver que efectivamente cumplen la condición de completitud, pues

$$\sum_{i=1}^2 E_i^\dagger E_i = |0\rangle\langle 0| + (1-p) |1\rangle\langle 1| + p|1\rangle\underbrace{\langle0|0\rangle}_{1}\langle1|= |0\rangle\langle0| + |1\rangle\langle1| = \mathbb{I}_A$$

Mientras que haciendo los productos en el orden inverso obtenemos

$$\sum_{i=1}^2 E_i E_i^\dagger = |0\rangle\langle 0| + (1-p) |1\rangle\langle 1| + p|0\rangle\underbrace{\langle1|1\rangle}_{1}\langle0|= (1-p)|1\rangle\langle1| + (1+p)|0\rangle\langle0| \neq \mathbb{I}_A.$$  

siempre que $p\neq 0$. Pues los operadores de Krauss no son necesariamente unitarios, y por tanto no necesariamente cumplen la condición de completitud en el orden inverso.

\textbf{b)} La evolución de $\rho_A$ está dada por $\rho_A' = \sum_{i=1}^2 E_i \rho_A E_i^\dagger$. Si tomamos $\rho_A$ como en el ejercicio anterior, entonces realizando la operación con el software de \textit{Mathematica}

\begin{align*}
    \rho_A' &= E_1 \rho_A E_1^\dagger + E_2 \rho_A E_2^\dagger \\
    &= \begin{pmatrix}
        1 & 0 \\
        0 & \sqrt{1-p}
    \end{pmatrix} \begin{pmatrix}
        \rho_{00} & \rho_{01} \\
        \rho_{10} & \rho_{11}
    \end{pmatrix} \begin{pmatrix}
        1 & 0 \\
        0 & \sqrt{1-p}
    \end{pmatrix} + 
    \begin{pmatrix}
        0 & \sqrt{p} \\
        0 & 0
    \end{pmatrix} 
    \begin{pmatrix}
        \rho_{00} & \rho_{01} \\
        \rho_{10} & \rho_{11}
    \end{pmatrix} 
    \begin{pmatrix}
        0 & 0\\
        \sqrt{p} & 0
    \end{pmatrix} \\[0.3cm]
    &= \begin{pmatrix}
        \rho_{00}+\rho_{11}p & \sqrt{1-p}\rho_{01} \\
        \sqrt{1-p}\rho_{10} & (1-p)\rho_{11}
    \end{pmatrix} 
\end{align*}


\end{exercise}


\begin{exercise}
    Consideramos un átomo de dos niveles con estados $|0_A\rangle$ con energía 0 y $|1_A\rangle$ con energía $\varepsilon$ y los estados de campo $|0_B\rangle$ y $|1_B\rangle=a^\dagger_w|0\rangle$ con $\hbar w$. Si suponemos que la evolución del sistema átomo + radiación es de la forma

\begin{align*}
    U = \begin{pmatrix*}
                1 & 0 & 0 & 0 \\
    0 & \cos(\theta) & \sin(\theta) & 0 \\
    0 & -\sin(\theta) & \cos(\theta) & 0 \\ 
    0 & 0 & 0 & 1       
    \end{pmatrix*}
\end{align*}

entonces, tomando la traza parcial de la evolución total

$$\rho'_A=\text{Tr}_C[U\rho_{A}\otimes |0_B\rangle\langle 0_B|U^\dagger]$$ 

obtenemos que los operadores de Krauss son

\end{exercise}















\end{document}
